\chapter{Software Requirements Specification}

The Software Requirements Specification (SRS) describes the requirements of
the proposed system and defines how the system is expected to function. It
specifies the system's main features, interfaces, inputs, outputs, and other
requirements needed for developing and evaluating the proposed
Custom RISC-V Instruction Extensions for Edge Video Anomaly Detection on the
Shakti C-Class processor.

\section{Introduction}

Edge video anomaly detection systems run deep learning models such as
MobileNetV2 with a GRU on small embedded processors placed close to the
camera. These processors execute the convolution and activation operations of
the model using general-purpose scalar instructions, where each
multiply-accumulate (MAC) operation is performed separately. On the baseline
Shakti C-Class core, one MAC takes 11 to 18 clock cycles, and a single
inference over eight video frames requires about 31.9 billion cycles. This
makes real-time processing on a low-cost edge processor difficult.

The proposed project addresses this problem by designing workload-driven
custom RISC-V instructions and integrating them into the open-source Shakti
C-Class core. Instead of adding instructions based on assumptions, the target
workload is first profiled on the processor itself to identify the operations
that consume the most execution time. Custom instructions are then designed
for these operations, implemented in the processor hardware, and called from C
programs through the RISC-V compiler. The current design includes three
instructions: FDOT2.S (packed dot product), PFMA.S (two-lane packed fused
multiply-add) and PRELU6.S (packed ReLU6 activation). Each instruction is
verified to produce results identical to standard RISC-V code, and its
performance is measured against the unmodified baseline core using the
processor's hardware cycle counters. This chapter provides the requirements
and functional structure needed to implement, test and evaluate the proposed
system.

\section{Overall Description}

The proposed system is a modified Shakti C-Class processor that executes deep
learning workloads for edge video anomaly detection faster than the original
core, using custom instructions designed from the measured behaviour of the
workload. The Shakti C-Class is a 64-bit, six-stage, in-order RISC-V core
implementing the RV64IMAFDC instruction set, written in Bluespec
SystemVerilog and simulated cycle by cycle using Verilator.

The development process begins with the target model and its layer
parameters. Each operation of the model, such as pointwise convolution,
depthwise convolution, the first convolution layer, the ReLU6 activation and
the GRU, is written as a bare-metal C kernel and executed on the baseline
processor. The workload profiling stage measures the cycles spent on each
operation and identifies the dominant kernels. Based on this profile, the
instruction design stage defines new instructions in the RISC-V custom-2
opcode space, packing two 32-bit floating-point values into one 64-bit
register so that a single instruction performs more useful work.

The hardware integration stage adds these instructions to the processor. The
decoder recognises the custom opcode and maps it onto the existing
floating-point path, and the floating-point unit is extended with a second
fused multiply-add unit and the logic required for the new operations. The
software stage calls the instructions from C code using the assembler
directive \texttt{.insn}. Finally, the verification and benchmarking stage
runs each kernel with standard code and with the custom instruction on the
simulated processor, checks that both outputs are bit-identical, reads the
cycle and instruction counts, and reports the speedup.

The overall design is modular. Each custom instruction is selected by its own
encoding and internal operation code, so instructions can be added, modified
or removed without changing the rest of the processor pipeline. The same
method can be applied to other models, such as the INT8 TinyMetro-Net model.

\section{External Interface Requirements}

The external interface of the proposed system defines how the user provides
the required inputs and receives the results. The user interacts with the
system through a command-line environment on a Linux system (Ubuntu under
WSL). The user provides the C benchmark program for the kernel to be
evaluated, and the system compiles it with the RISC-V GCC toolchain, converts
it into a memory image using \texttt{elf2hex}, and loads it into the
simulated processor. A single command can be used to build, run and report
the results of a selected benchmark.

The hardware interface is the Shakti C-Class processor model generated from
the Bluespec source code. The Bluespec compiler converts the modified source
into Verilog, and Verilator builds a cycle-accurate simulator from it. The
software interface between programs and the new hardware is the RISC-V
instruction set itself: the custom instructions use the reserved custom-2
opcode (\texttt{0x5B}) with the R4 instruction layout and are emitted from C
code through the \texttt{.insn} directive, so no change to the compiler is
required.

During execution, the simulated processor writes an instruction trace
(\texttt{rtl.dump}), records the cycle and instruction counts using the
\texttt{mcycle} and \texttt{minstret} hardware counters, and reports pass or
fail through the \texttt{tohost} memory location. The system presents the
final result to the user as the number of cycles taken by the baseline and
custom versions, the resulting speedup, and whether the outputs of both
versions are identical. The system uses only locally stored source files,
benchmark programs, configuration files and result files. No network
connection is required for building, simulating or evaluating the processor.

\section{System Features}

The proposed system consists of several features that work together to
accelerate the target workload. The process begins with workload profiling,
where every operation of the MobileNetV2 + GRU model is executed on the
baseline processor with its real layer dimensions and timed using the
hardware cycle counter. The measured cost of each operation is combined with
the number of times it runs in one inference to identify the operations that
dominate the execution time.

The identified operations are then addressed through custom instruction
design. FDOT2.S computes a two-element dot product, PFMA.S performs two fused
multiply-add operations in parallel on two packed values, and PRELU6.S applies
the ReLU6 activation to two packed values in a single cycle. These
instructions are integrated into the processor through decoder changes that
map the custom opcode onto the existing floating-point path, and through
floating-point unit changes that add a second fused multiply-add unit, the
two-step dot product logic, result packing and the ReLU6 logic.

The instructions are used by C kernels for pointwise convolution, depthwise
convolution, the first convolution layer and ReLU6. The verification feature
runs each kernel in both standard and custom forms, compares all outputs bit
for bit, and runs the official RISC-V \texttt{fmadd} test to confirm that
standard instructions still behave correctly. Finally, the benchmarking
feature reports the cycles per operation and the speedup of each kernel, and
estimates the execution time of a complete inference on the modified
processor compared with the baseline.

\section{Other Non-Functional Requirements}

The proposed system should produce results that are bit-identical to standard
RISC-V code for the same inputs, and the measured cycle counts should be
repeatable when the same program and processor build are used. The custom
instructions should provide a measurable improvement, with a target of at
least $1.5\times$ speedup on the dominant convolution kernels, while keeping
the hardware additions small. The existing RV64IMAFDC instructions must
continue to function correctly after every change, which is confirmed using
the official riscv-tests.

The system should be easy to reproduce, using only open-source tools: the
Bluespec compiler, Verilator and the RISC-V GCC toolchain. Building,
running and evaluating a benchmark should be possible through simple
commands, and the results should clearly show the baseline cycles, the custom
cycles, the speedup and the pass or fail status. The design should also
respect the constraints of an edge processor, including its limited memory
such as the 16\,KB data cache.

The design should be modular and maintainable, with each custom instruction
implemented as a separate, clearly identified part of the decoder and
floating-point unit, and with the design decisions documented. This allows
individual instructions to be tested, modified or replaced without affecting
the rest of the processor. The system should also be extensible, so that
additional instructions, such as INT8 instructions for the TinyMetro-Net
model or a pipelined floating-point unit, can be added in the future while the
existing instructions continue to work.
